\section{Length}
\label{section-length}

\begin{definition}
\label{definition-length}
Let $R$ be a ring. For any $R$-module $M$
we define the {\it length} of $M$ over $R$ by the
formula
$$
\text{length}_R(M)
=
\sup
\{
n
\mid
\exists\ 0 = M_0 \subset M_1 \subset \ldots \subset M_n = M,
\text{ }M_i \not = M_{i + 1}
\}.
$$
\end{definition}

\noindent
In other words it is the supremum of the lengths of chains
of submodules. There is an obvious notion of when a chain
of submodules is a refinement of another. This gives a
partial ordering on the collection of all chains of submodules,
with the smallest chain having the shape $0 = M_0 \subset M_1 = M$
if $M$ is not zero.
We note the obvious fact that if the length of
$M$ is finite, then every chain can be refined to a
maximal chain. But it is not as obvious that all maximal
chains have the same length (as we will see later).

\begin{lemma}
\label{lemma-finite-length-finite}
\begin{slogan}
Modules of finite length are finite.
\end{slogan}
Let $R$ be a ring.
Let $M$ be an $R$-module.
If $\text{length}_R(M) < \infty$ then $M$ is a finite $R$-module.
\end{lemma}

\begin{proof}
Omitted.
\end{proof}

\begin{lemma}
\label{lemma-length-additive}
\begin{slogan}
Length is additive in short exact sequences.
\end{slogan}
If $0 \to M' \to M \to M'' \to 0$
is a short exact sequence of modules over $R$ then
the length of $M$ is the sum of the
lengths of $M'$ and $M''$.
\end{lemma}

\begin{proof}
Given filtrations of $M'$ and $M''$ of lengths $n', n''$
it is easy to make a corresponding filtration of $M$
of length $n' + n''$. Thus we see that $\text{length}_R M
\geq \text{length}_R M' + \text{length}_R M''$.
Conversely, given a filtration
$M_0 \subset M_1 \subset \ldots \subset M_n$ of
$M$ consider the induced filtrations
$M_i' = M_i \cap M'$ and $M_i'' = \Im(M_i \to M'')$.
Let $n'$ (resp.\ $n''$) be the number of steps in the filtration
$\{M'_i\}$ (resp.\ $\{M''_i\}$).
If $M_i' = M_{i + 1}'$ and $M_i'' = M_{i + 1}''$ then
$M_i = M_{i + 1}$. Hence we conclude that $n' + n'' \geq n$.
Combined with the earlier result we win.
\end{proof}

\begin{lemma}
\label{lemma-length-infinite}
Let $R$ be a local ring with maximal ideal $\mathfrak m$.
If $M$ is an $R$-module and $\mathfrak m^n M \not = 0$ for all $n \geq 0$,
then $\text{length}_R(M) = \infty$. In other words, if $M$
has finite length then $\mathfrak m^nM = 0$ for some $n$.
\end{lemma}

\begin{proof}
Assume $\mathfrak m^n M \not = 0$ for all $n\geq 0$.
Choose $x \in M$ and $f_1, \ldots, f_n \in \mathfrak m$
such that $f_1f_2 \ldots f_n x \not = 0$.
The first $n$ steps in the filtration
$$
0 \subset R f_1 \ldots f_n x \subset R f_1 \ldots f_{n - 1} x
\subset \ldots \subset R x \subset M
$$
are distinct. For example, if
$R f_1 x = R f_1 f_2 x$ , then $f_1 x = g f_1 f_2 x$ for some $g$,
hence $(1 - gf_2) f_1 x = 0$ hence $f_1 x = 0$ as $1 - gf_2$ is a unit
which is a contradiction with the choice of $x$ and $f_1, \ldots, f_n$.
Hence the length is infinite.
\end{proof}

\begin{lemma}
\label{lemma-length-independent}
Let $R \to S$ be a ring map. Let $M$ be an $S$-module.
We always have $\text{length}_R(M) \geq \text{length}_S(M)$.
If $R \to S$ is surjective then equality holds.
\end{lemma}

\begin{proof}
A filtration of $M$ by $S$-submodules gives rise a filtration
of $M$ by $R$-submodules. This proves the inequality.
And if $R \to S$ is surjective, then any $R$-submodule
of $M$ is automatically an $S$-submodule. Hence equality
in this case.
\end{proof}

\begin{lemma}
\label{lemma-dimension-is-length}
Let $R$ be a ring with maximal ideal $\mathfrak m$.
Suppose that $M$ is an $R$-module with
$\mathfrak m M  =  0$. Then the length of $M$ as
an $R$-module agrees with the dimension of $M$ as
a $R/\mathfrak m$ vector space.
The length is finite if and only if $M$ is a finite $R$-module.
\end{lemma}

\begin{proof}
The first part is a special case of Lemma \ref{lemma-length-independent}.
Thus the length is finite if and only if $M$ has a finite basis
as a $R/\mathfrak m$-vector space if and only if $M$ has a finite
set of generators as an $R$-module.
\end{proof}

\begin{lemma}
\label{lemma-length-localize}
Let $R$ be a ring. Let $M$ be an $R$-module. Let $S \subset R$ be
a multiplicative subset. Then
$\text{length}_R(M) \geq \text{length}_{S^{-1}R}(S^{-1}M)$.
\end{lemma}

\begin{proof}
Any submodule $N' \subset S^{-1}M$ is of the form
$S^{-1}N$ for some $R$-submodule $N \subset M$, by Lemma
\ref{lemma-submodule-localization}. The lemma follows.
\end{proof}

\begin{lemma}
\label{lemma-length-finite}
Let $R$ be a ring with finitely generated
maximal ideal $\mathfrak m$. (For example $R$ Noetherian.)
Suppose that $M$ is a finite $R$-module with
$\mathfrak m^n M  =  0$ for some $n$.
Then $\text{length}_R(M) < \infty$.
\end{lemma}

\begin{proof}
Consider the filtration
$0 = \mathfrak m^n M \subset
\mathfrak m^{n-1} M \subset
\ldots \subset \mathfrak m M \subset M$.
All of the subquotients are finitely generated $R$-modules
to which Lemma \ref{lemma-dimension-is-length} applies. We conclude
by additivity, see Lemma \ref{lemma-length-additive}.
\end{proof}

\begin{definition}
\label{definition-simple-module}
Let $R$ be a ring. Let $M$ be an $R$-module.
We say $M$ is {\it simple} if $M \not = 0$ and
every submodule of $M$ is either equal to $M$ or
to $0$.
\end{definition}

\begin{lemma}
\label{lemma-characterize-length-1}
Let $R$ be a ring. Let $M$ be an $R$-module.
The following are equivalent:
\begin{enumerate}
\item $M$ is simple,
\item $\text{length}_R(M) = 1$, and
\item $M \cong R/\mathfrak m$ for some maximal ideal
$\mathfrak m \subset R$.
\end{enumerate}
\end{lemma}

\begin{proof}
Let $\mathfrak m$ be a maximal ideal of $R$.
By Lemma \ref{lemma-dimension-is-length} the module
$R/\mathfrak m$ has length $1$. The equivalence of
the first two assertions is tautological.
Suppose that $M$ is simple. Choose $x \in M$, $x \not = 0$.
As $M$ is simple we have $M = R \cdot x$.
Let $I \subset R$ be the annihilator of $x$, i.e.,
$I = \{f \in R \mid fx = 0\}$. The map $R/I \to M$,
$f \bmod I \mapsto fx$ is an isomorphism, hence
$R/I$ is a simple $R$-module. Since $R/I \not = 0$ we see $I \not = R$.
Let $I \subset \mathfrak m$ be a maximal ideal containing $I$.
If $I \not = \mathfrak m$, then $\mathfrak m /I \subset R/I$
is a nontrivial submodule contradicting the simplicity
of $R/I$. Hence we see $I = \mathfrak m$ as desired.
\end{proof}

\begin{lemma}
\label{lemma-simple-pieces}
Let $R$ be a ring. Let $M$ be a finite length $R$-module.
Choose any maximal chain of submodules
$$
0 = M_0 \subset M_1 \subset M_2 \subset \ldots \subset M_n = M
$$
with $M_i \not = M_{i-1}$, $i = 1, \ldots, n$. Then
\begin{enumerate}
\item $n = \text{length}_R(M)$,
\item each $M_i/M_{i-1}$ is simple,
\item each $M_i/M_{i-1}$ is of the form
$R/\mathfrak m_i$ for some maximal ideal $\mathfrak m_i$,
\item given a maximal ideal $\mathfrak m \subset R$
we have
$$
\# \{i \mid \mathfrak m_i = \mathfrak m\}
=
\text{length}_{R_{\mathfrak m}} (M_{\mathfrak m}).
$$
\end{enumerate}
\end{lemma}

\begin{proof}
If $M_i/M_{i-1}$ is not simple then we can refine the filtration
and the filtration is not maximal. Thus we see that $M_i/M_{i-1}$
is simple. By Lemma \ref{lemma-characterize-length-1} the modules
$M_i/M_{i-1}$ have length $1$ and are of the form $R/\mathfrak m_i$
for some maximal ideals $\mathfrak m_i$. By additivity of length,
Lemma \ref{lemma-length-additive}, we see $n = \text{length}_R(M)$.
Since localization is exact, we see that
$$
0 = (M_0)_{\mathfrak m}
\subset (M_1)_{\mathfrak m}
\subset (M_2)_{\mathfrak m}
\subset \ldots
\subset (M_n)_{\mathfrak m} = M_{\mathfrak m}
$$
is a filtration of $M_{\mathfrak m}$ with successive quotients
$(M_i/M_{i-1})_{\mathfrak m}$. Thus the last statement follows
directly from the fact that given maximal ideals $\mathfrak m$,
$\mathfrak m'$ of $R$ we have
$$
(R/\mathfrak m')_{\mathfrak m}
\cong
\left\{
\begin{matrix}
0 &
\text{if } \mathfrak m \not = \mathfrak m', \\
R_{\mathfrak m}/\mathfrak m R_{\mathfrak m} &
\text{if } \mathfrak m  = \mathfrak m'
\end{matrix}
\right.
$$
This we leave to the reader.
\end{proof}

\begin{lemma}
\label{lemma-pushdown-module}
Let $A$ be a local ring with maximal ideal $\mathfrak m$.
Let $B$ be a semi-local ring with maximal ideals $\mathfrak m_i$,
$i = 1, \ldots, n$.
Suppose that $A \to B$ is a homomorphism such that each $\mathfrak m_i$
lies over $\mathfrak m$ and such that
$$
[\kappa(\mathfrak m_i) : \kappa(\mathfrak m)] < \infty.
$$
Let $M$ be a $B$-module of finite length.
Then
$$
\text{length}_A(M) = \sum\nolimits_{i = 1, \ldots, n}
[\kappa(\mathfrak m_i) : \kappa(\mathfrak m)]
\text{length}_{B_{\mathfrak m_i}}(M_{\mathfrak m_i}),
$$
in particular $\text{length}_A(M) < \infty$.
\end{lemma}

\begin{proof}
Choose a maximal chain
$$
0 = M_0
\subset M_1
\subset M_2
\subset \ldots
\subset M_m = M
$$
by $B$-submodules as in Lemma \ref{lemma-simple-pieces}.
Then each quotient $M_j/M_{j - 1}$ is isomorphic to
$\kappa(\mathfrak m_{i(j)})$ for some $i(j) \in \{1, \ldots, n\}$.
Moreover
$\text{length}_A(\kappa(\mathfrak m_i)) =
[\kappa(\mathfrak m_i) : \kappa(\mathfrak m)]$ by
Lemma \ref{lemma-dimension-is-length}. The lemma follows
by additivity of lengths (Lemma \ref{lemma-length-additive}).
\end{proof}

\begin{lemma}
\label{lemma-pullback-module}
Let $A \to B$ be a flat local homomorphism of local rings.
Then for any $A$-module $M$ we have
$$
\text{length}_A(M) \text{length}_B(B/\mathfrak m_AB)
=
\text{length}_B(M \otimes_A B).
$$
In particular, if $\text{length}_B(B/\mathfrak m_AB) < \infty$
then $M$ has finite length if and only if $M \otimes_A B$ has finite length.
\end{lemma}

\begin{proof}
The ring map $A \to B$ is faithfully flat by
Lemma \ref{lemma-local-flat-ff}.
Hence if $0 = M_0 \subset M_1 \subset \ldots \subset M_n = M$
is a chain of length $n$ in $M$, then the corresponding chain
$0 = M_0 \otimes_A B \subset M_1 \otimes_A B \subset
\ldots \subset M_n \otimes_A B = M \otimes_A B$ has length $n$
also. This proves
$\text{length}_A(M) = \infty \Rightarrow
\text{length}_B(M \otimes_A B) = \infty$.
Next, assume $\text{length}_A(M) < \infty$. In this case we see
that $M$ has a filtration of length $\ell = \text{length}_A(M)$
whose quotients are $A/\mathfrak m_A$. Arguing as above
we see that $M \otimes_A B$ has a filtration of length $\ell$
whose quotients are isomorphic to
$B \otimes_A A/\mathfrak m_A =  B/\mathfrak m_AB$.
Thus the lemma follows.
\end{proof}

\begin{lemma}
\label{lemma-pullback-transitive}
Let $A \to B \to C$ be flat local homomorphisms of local rings. Then
$$
\text{length}_B(B/\mathfrak m_A B)
\text{length}_C(C/\mathfrak m_B C)
=
\text{length}_C(C/\mathfrak m_A C)
$$
\end{lemma}

\begin{proof}
Follows from Lemma \ref{lemma-pullback-module} applied to the ring map
$B \to C$ and the $B$-module $M = B/\mathfrak m_A B$
\end{proof}











